\chapter{Datenschutz und Privatsphäre bei KI-Tools}\label{ch:K06_Datenschutz}

\begin{myremark}
Wenn Sie ein KI-Tool nutzen, teilen Sie Daten mit einem Unternehmen. In diesem
Kapitel analysieren Sie, welche Daten KI-Anbieter sammeln, was die AGBs
(Allgemeinen Geschäftsbedingungen) dazu sagen, und welche datenschutzkonformen
Alternativen für den Schulbereich existieren.
\end{myremark}

% ─────────────────────────────────────────────────────────────────────────────
\section{Welche Daten sammeln KI-Anbieter?}

\begin{mysummary}
Typische Daten, die KI-Plattformen sammeln:
\begin{itemize}
    \item \textbf{Prompt-Inhalte}: Alles, was Sie eingeben, wird auf den Servern
          des Anbieters gespeichert.
    \item \textbf{Metadaten}: IP-Adresse, Zeitstempel, Geräteinformationen.
    \item \textbf{Nutzungsverhalten}: Wie oft, wie lange und welche Funktionen
          Sie nutzen.
    \item \textbf{Konversationsverläufe}: Alle bisherigen Gespräche
          (sofern nicht deaktiviert).
\end{itemize}
\end{mysummary}

\begin{myattention}
\textbf{Grundregel}: Geben Sie \emph{niemals} persönliche Daten (Name, Adresse,
AHV-Nummer, Passwörter, Bankdaten) in ein KI-Tool ein. Auch scheinbar harmlose
Informationen können in Kombination problematisch sein.
\end{myattention}

% ─────────────────────────────────────────────────────────────────────────────
\section{AGB-Analyse: Was erlauben Sie dem Anbieter?}

\begin{myprocedure}[Ablauf der AGB-Analyse]
\begin{enumerate}
    \item Öffnen Sie die AGBs des zugewiesenen KI-Tools (Links in \cref{sec:agb-links}).
    \item Suchen Sie die relevanten Abschnitte zu Datenspeicherung, Nutzung
          der Eingaben und Weitergabe an Dritte.
    \item Beantworten Sie die Leitfragen.
    \item Vergleichen Sie Ihre Ergebnisse mit einer anderen Gruppe.
\end{enumerate}
\end{myprocedure}

\subsection{Links zu den AGBs}\label{sec:agb-links}

\begin{itemize}
    \item \textbf{ChatGPT (OpenAI)}: \url{https://openai.com/policies/terms-of-use}
    \item \textbf{Claude (Anthropic)}: \url{https://www.anthropic.com/legal/consumer-terms}
    \item \textbf{Gemini (Google)}: \url{https://policies.google.com/terms}
    \item \textbf{Fobizz}: \url{https://fobizz.com/agb}
    \item \textbf{Brian}: (Schulinterne Lizenz – fragen Sie Ihre Lehrperson nach
          den Datenschutzbestimmungen)
\end{itemize}

\begin{myexercise}[– AGB-Analyse]
Analysieren Sie die AGBs \textbf{eines} der oben genannten Dienste
(Ihre Lehrperson weist Ihnen einen zu). Beantworten Sie folgende Fragen:

\begin{enumerate}
    \item \textbf{Datenspeicherung}: Wie lange werden Ihre Daten gespeichert?
          Wo (in welchem Land)?
    \item \textbf{Nutzung der Eingaben}: Darf der Anbieter Ihre Prompts für
          das Training neuer Modelle nutzen? Können Sie das deaktivieren?
    \item \textbf{Datenweitergabe}: Werden Ihre Daten an Dritte weitergegeben?
          Unter welchen Umständen?
    \item \textbf{Ihre Rechte}: Haben Sie das Recht auf Löschung Ihrer Daten?
          Wie können Sie dieses Recht ausüben?
    \item \textbf{Anwendbares Recht}: Welches Recht gilt bei Streitigkeiten
          (Schweizer Recht? US-amerikanisches Recht?)?
\end{enumerate}

Halten Sie Ihre Antworten in der folgenden Tabelle fest:

\begin{center}
\begin{tabular}{p{4cm}p{9cm}}
    \toprule
    \textbf{Frage}               & \textbf{Antwort (mit Zitatstelle aus den AGBs)} \\
    \midrule
    Datenspeicherungsdauer       &                                                  \\[1cm]
    Speicherort (Land)           &                                                  \\[1cm]
    Training mit Prompts         &                                                  \\[1cm]
    Datenweitergabe              &                                                  \\[1cm]
    Löschungsrecht               &                                                  \\[1cm]
    Anwendbares Recht            &                                                  \\[1cm]
    \bottomrule
\end{tabular}
\end{center}
\end{myexercise}

% ─────────────────────────────────────────────────────────────────────────────
\section{Multiple-Choice: Datenschutz und KI}

\begin{myexercise}[– MC-Quiz Datenschutz]
Beantworten Sie die folgenden Multiple-Choice-Fragen. Mehrere Antworten
können richtig sein.

\begin{enumerate}

\item Was sollte man \textbf{niemals} in ein öffentliches KI-Tool eingeben?
\begin{aenum}
    \item Eine Frage über den Klimawandel
    \item Den eigenen Vor- und Nachnamen kombiniert mit der Wohnadresse
    \item Eine Zusammenfassung des Klassenausflugs
    \item Das eigene Passwort für ein Online-Konto
    \item Einen Text auf Englisch übersetzen lassen
\end{aenum}

\item Welche Aussage über die DSGVO (Datenschutz-Grundverordnung) ist korrekt?
\begin{aenum}
    \item Die DSGVO gilt weltweit für alle Unternehmen.
    \item Die DSGVO schützt die persönlichen Daten von Personen in der EU.
    \item Wenn ein US-Unternehmen Daten von EU-Bürgerinnen verarbeitet,
          muss es die DSGVO einhalten.
    \item In der Schweiz gilt das DSG (Datenschutzgesetz) als äquivalentes Gesetz.
    \item DSGVO-konformes Speichern bedeutet, dass Daten niemals gespeichert werden.
\end{aenum}

\item Was bedeutet es, wenn ein KI-Anbieter \enquote{opt-out from training}
      anbietet?
\begin{aenum}
    \item Ihre Daten werden vollständig gelöscht.
    \item Ihre Konversationen werden nicht für das Training neuer Modelle verwendet.
    \item Sie können die KI nicht mehr nutzen.
    \item Der Anbieter löscht alle Metadaten.
    \item Ihre Prompts werden weiterhin gespeichert, aber nicht für Training genutzt.
\end{aenum}

\item Welche der folgenden Tools sind für den Schulbereich in der Schweiz
      datenschutzkonforme Alternativen?
\begin{aenum}
    \item ChatGPT (kostenlose Version)
    \item Fobizz (mit Schullizenz)
    \item Brian (mit Schullizenz)
    \item Gemini (Google-Account erforderlich)
    \item Microsoft Copilot (mit schulischem Microsoft-365-Account)
\end{aenum}

\item Ein Schüler gibt folgenden Prompt ein: \enquote{Ich heisse Max Muster,
      gehe in die Klasse 3a am Gymnasium Musterstadt und brauche Hilfe bei
      meiner Mathe-Hausaufgabe.} Was ist problematisch?
\begin{aenum}
    \item Gar nichts – es geht nur um Mathe.
    \item Der vollständige Name in Kombination mit Schule und Klasse
          sind persönliche Daten.
    \item KI-Tools dürfen keine Mathe-Hausaufgaben bearbeiten.
    \item Der Anbieter könnte diese Daten mit anderen Quellen kombinieren.
    \item Der Name allein reicht, um die Person zu identifizieren.
\end{aenum}

\end{enumerate}
\end{myexercise}

% ─────────────────────────────────────────────────────────────────────────────
\section{Datenschutzkonforme KI-Tools für die Schule}

\subsection{Fobizz}

\begin{myexample}
\textbf{Fobizz} (\url{https://fobizz.com}) ist eine deutsche Bildungsplattform,
die DSGVO-konforme KI-Tools für Schulen anbietet. Features:
\begin{itemize}
    \item KI-Assistent für Schülerinnen und Schüler (basierend auf GPT-4).
    \item Keine Weitergabe von Daten an Dritte.
    \item Server in der EU.
    \item Werkzeuge für Lehrpersonen: Unterrichtsvorbereitung, Quiz-Erstellung etc.
\end{itemize}
\end{myexample}

\subsection{Brian}

\begin{myexample}
\textbf{Brian} ist ein schweizspezifisches KI-Tool für Schulen, das besondere
Datenschutzanforderungen des Schweizer Bildungssystems berücksichtigt.
Fragen Sie Ihre Lehrperson nach den genauen Datenschutzbestimmungen
und dem Zugangsweg für Ihre Schule.
\end{myexample}

\begin{myexercise}[– Vergleich: Datenschutz Fobizz vs. ChatGPT]
Vergleichen Sie Fobizz und ChatGPT (kostenlose Version) anhand der
folgenden Kriterien und füllen Sie die Tabelle aus:

\begin{center}
\begin{tabular}{p{4cm}p{4.5cm}p{4.5cm}}
    \toprule
    \textbf{Kriterium}             & \textbf{Fobizz}  & \textbf{ChatGPT (kostenlos)} \\
    \midrule
    Serverstandort                 &                  &                              \\[0.5cm]
    Daten für Training genutzt?    &                  &                              \\[0.5cm]
    Anwendbares Recht              &                  &                              \\[0.5cm]
    Altersbeschränkung             &                  &                              \\[0.5cm]
    Schullizenz vorhanden?         &                  &                              \\[0.5cm]
    Eingeschränkter Funktionsumfang?&                 &                              \\[0.5cm]
    \bottomrule
\end{tabular}
\end{center}
\end{myexercise}

% ─────────────────────────────────────────────────────────────────────────────
\section{Exkurs: Das Schweizer Datenschutzgesetz (DSG)}

\begin{myexcursion}
Das revidierte \textbf{Datenschutzgesetz (DSG)}, in der Schweiz seit September 2023
in Kraft, schützt die Persönlichkeit und Grundrechte von Menschen,
über die Personendaten bearbeitet werden. Wesentliche Rechte:

\begin{itemize}
    \item \textbf{Auskunftsrecht}: Sie dürfen fragen, welche Daten über Sie
          gespeichert sind.
    \item \textbf{Berichtigungsrecht}: Sie dürfen falsche Daten korrigieren lassen.
    \item \textbf{Löschungsrecht}: Unter bestimmten Bedingungen können Sie
          die Löschung Ihrer Daten verlangen.
    \item \textbf{Recht auf Datenübertragbarkeit}: Ihre Daten müssen Ihnen
          in einem maschinenlesbaren Format herausgegeben werden.
\end{itemize}

Das DSG orientiert sich stark an der europäischen DSGVO und ist mit ihr
weitgehend kompatibel.
\end{myexcursion}

\begin{mychallenge}[Datenschutzfolgenabschätzung]
\begin{myattention}
\textbf{Optional}: Diese Aufgabe eignet sich als vertiefte Hausaufgabe oder
Gruppenarbeit.
\end{myattention}

Stellen Sie sich vor, Ihre Schule möchte ein neues KI-Tool einführen, das
Hausaufgaben korrigiert. Erstellen Sie eine vereinfachte
\textbf{Datenschutzfolgenabschätzung} (DPIA – Data Protection Impact Assessment)
mit folgenden Punkten:

\begin{enumerate}
    \item Welche Daten werden verarbeitet?
    \item Wer hat Zugriff auf die Daten?
    \item Welche Risiken entstehen für die Schülerinnen und Schüler?
    \item Welche technischen und organisatorischen Massnahmen könnten
          die Risiken reduzieren?
    \item Empfehlen Sie das Tool einzuführen oder nicht? Begründen Sie.
\end{enumerate}
\end{mychallenge}
